\documentclass[11pt]{amsart}

\usepackage[T1]{fontenc}
\usepackage{lmodern}
\usepackage{microtype}
\usepackage{mathtools,amssymb,amsthm}
\usepackage[hidelinks]{hyperref}

\hypersetup{
  pdftitle={A mod 31 squared counterexample to a generalized cubic partition
            congruence of Das, Maity and Saikia}
}

\newtheorem{theorem}{Theorem}
\theoremstyle{remark}
\newtheorem{remark}[theorem]{Remark}

\DeclareMathOperator{\vv}{v}

\emergencystretch=3em

\title[A mod $31^2$ counterexample to a cubic partition congruence]
{A mod $31^2$ counterexample to a generalized cubic partition congruence
of Das, Maity and Saikia}

\date{}

\subjclass[2020]{11P83, 05A17, 11F03}
\keywords{generalized cubic partition, partition congruence, $p$-adic
valuation, Ramanujan--Kolberg, counterexample}

\begin{document}

\begin{abstract}
For $c\ge1$ let $a_c(n)$ be the number of partitions of $n$ in which the even
parts may take $c$ colours, so that $\sum_{n\ge0}a_c(n)q^n=1/(f_1f_2^{c-1})$
with $f_k=(q^k;q^k)_\infty$. In the Concluding Remarks of the preprint
\cite{DasMaitySaikia}, Das, Maity and Saikia conjecture three congruences to a
prime-square modulus, the first of which is
$a_{25}(31^2n+644)\equiv0\pmod{31^2}$ for all $n\ge0$; they state explicitly
that they leave it open because the Ramanujan--Kolberg bound is too large for
their computer. We refute it at $n=2$. The exact value of $a_{25}(2566)$ is a
$182$-digit integer, printed in full below; it satisfies
$a_{25}(2566)\equiv682=22\cdot31\pmod{31^2}$, so its $31$-adic valuation is
exactly $1$ and the congruence fails. What survives is
the first power: $a_{25}(31^2n+644)\equiv0\pmod{31}$ does hold for all $n$, but
that is a published theorem of Guadalupe \cite{Guadalupe}, and it is only its
lift to $31^2$ that is false.
\end{abstract}

\maketitle

\section{The conjecture}\label{sec:conj}

Write $f_k=(q^k;q^k)_\infty$. Following Amdeberhan, Sellers and Singh
\cite{ASS}, a \emph{generalized cubic partition} of $n$ is a partition of $n$ in
which each \emph{even} part may appear in one of $c\ge1$ colours, and $a_c(n)$
denotes the number of these, with $a_c(0)=1$. Equivalently, and this is the
definition used throughout,
\begin{equation}\label{eq:gf}
  \sum_{n\ge0}a_c(n)q^n \;=\; \frac{1}{f_1 f_2^{\,c-1}}.
\end{equation}
Thus $a_1(n)$ is the ordinary partition function (OEIS A000041) and $a_2(n)$
counts Chan's cubic partitions (OEIS A002513) \cite{OEIS}.

\begin{remark}\label{rem:defn}
The source \cite{DasMaitySaikia} defines $a_c(n)$ by \eqref{eq:gf}, but its
prose reads ``each part may appear in $c\ge1$ different colors'', which would
give $1/f_1^{\,c}$ instead. The word \emph{even} was dropped; both \cite{ASS},
where the function originates, and \cite{Guadalupe} carry it. Equation
\eqref{eq:gf} is the authoritative object, and nothing below is sensitive to the
choice: under the prose reading $1/f_1^{25}$ the congruence at issue already
fails at $n=0$, with residue $550$ rather than $0$ modulo $961$.
\end{remark}

The statement refuted here is the first displayed congruence of the first
conjecture in the Concluding Remarks of \texttt{arXiv:2503.19399v2}
\cite{DasMaitySaikia}, quoted below from the source of that version. We identify
it by that position and not by a number: no claim is made here about how the
source numbers this conjecture, and \cite{GhoshalJana} refers to its third
congruence as ``Conjecture 7.3''.

\medskip
\noindent\textbf{Conjecture} (Concluding Remarks of
\texttt{arXiv:2503.19399v2} \cite{DasMaitySaikia})\textbf{.}
\emph{For all $n\ge0$, we have}
\begin{align}
  a_{25}(31^2n+644)&\equiv0 \pmod{31^2},\label{eq:line1}\\
  a_{41}(47^2n+256)&\equiv0 \pmod{47^2},\label{eq:line2}\\
  a_{61}(67^2n+555)&\equiv0 \pmod{67^2}.\label{eq:line3}
\end{align}

\medskip
\noindent
The single quantifier ``for all $n\ge0$'' scopes the whole display: there is no
side hypothesis, no exceptional set and no ``sufficiently large $n$'', and the
constants $c=25$, $p=31$ and $r=644$ are fixed, so one explicit $n$ refutes
\eqref{eq:line1}. The authors print Ramanujan--Kolberg parameters for each of
the three lines, including the bound $\lfloor\nu\rfloor=3813$ of Radu's
algorithm \cite{Radu} for \eqref{eq:line1}, and then write, verbatim:
\begin{quote}
``Since the bound $\nu$ is large for the computation to take a significant
amount of time in a modest computer, we leave the above as an open problem.''
\end{quote}
Deciding \eqref{eq:line1} at a single $n$, however, needs one coefficient of one
truncated series and not the cells of a certificate.

\section{The counterexample}\label{sec:refute}

\begin{theorem}\label{thm:main}
Congruence \eqref{eq:line1}, the first displayed congruence of the conjecture
quoted above from \texttt{arXiv:2503.19399v2} \cite{DasMaitySaikia}, is false.
Explicitly it
fails at $n=2$: with
$N=31^2\cdot2+644=1922+644=2566$,
\[
  a_{25}(2566)\;\equiv\;682\;=\;22\cdot31 \pmod{31^2},
  \qquad\text{so}\qquad \vv_{31}\bigl(a_{25}(2566)\bigr)=1 .
\]
In particular $31\mid a_{25}(2566)$ while $31^2\nmid a_{25}(2566)$.
\end{theorem}

\begin{proof}
Let $W$ be the $182$-digit integer
\begin{center}\ttfamily\footnotesize
1471465862381281557802668516631902844710 \\
2593034654589643689045403133888112437860 \\
2701871704379341548682799083362285352501 \\
0333330703396635284385120674666433997413 \\
0010853309981128071099
\end{center}
and let $Q$ be the $180$-digit integer
\begin{center}\ttfamily\footnotesize
4746664072197682444524737150425493047452 \\
4493660176095624803372268173832620767291 \\
1941521627030134028009029301168662427422 \\
6881711946440758981887486047311077411009 \\
68092042257455744229
\end{center}
Then $a_{25}(2566)=W$, and $31Q=W$, so $31\mid W$. Reducing $Q$ modulo $31$
gives $22\ne0$, whence $\vv_{31}(W)=1$ exactly and
$W\equiv 31\cdot 22=682\pmod{961}$. Since $682\ne0$ and
$N=2566=31^2\cdot2+644$ lies on the progression of \eqref{eq:line1} with
$n=2\ge0$, the congruence asserted there fails at that cell.

Only the identity $a_{25}(2566)=W$ needs a computation, and it is verified in
\S\ref{sec:verify}; the rest is the division of one printed integer by another.
A reader with any computer algebra system needs the coefficient of $q^{2566}$ in
$1/\bigl((q;q)_\infty(q^2;q^2)_\infty^{24}\bigr)$, a $2567$-term truncation, and
one division by $31$.
\end{proof}

\section{The companion congruence modulo $31$}\label{sec:why}

That the whole progression is divisible by the first power of $31$ is not new
and is not claimed here. Theorem 3 of Guadalupe \cite{Guadalupe} reads: \emph{let
$p\ge7$ be a prime with $p\equiv3,7\pmod 8$; then for all $n\ge0$,
$a_{p-6}\bigl(pn+13(p^2-1)/24\bigr)\equiv0\pmod p$.} At $p=31$ we have
$31\equiv7\pmod 8$ and $13(31^2-1)/24=520=31\cdot16+24$, while
$961n+644=31(31n+20)+24$ with $31n+20\ge20\ge16$ for every $n\ge0$. Hence
\begin{equation}\label{eq:published}
  a_{25}(31^2n+644)\equiv0 \pmod{31}\qquad\text{for all }n\ge0
\end{equation}
is a published theorem, predating \cite{DasMaitySaikia} and cited in it. What
\eqref{eq:line1} adds to \eqref{eq:published} is the second power, and it is
only that lift which Theorem~\ref{thm:main} refutes.

\section{The computation and \texttt{verify.py}}\label{sec:verify}

The single coefficient that decides Theorem~\ref{thm:main} is the $2567$th of a
truncated $q$-series, so the computation is small; what needs care is that the
right series is being expanded. The program \texttt{verify.py} accompanying this
note is Python 3.9 or later, standard library only, exact integer arithmetic, no
external data file. It recomputes $a_{25}(0),\dots,a_{25}(2566)$ over
$\mathbb{Z}$ by two routes that share no arithmetic ---
(i) building $D=f_1f_2^{24}$ from the generalized pentagonal numbers and
inverting it by $A_0=1$, $A_n=-\sum_{j=1}^{n}D_jA_{n-j}$, in which no division
and no modular inverse occurs because $D_0=1$; and
(ii) the logarithmic-derivative recurrence
$n\,a_c(n)=\sum_{j=1}^{n}B(j)a_c(n-j)$ with
$B(j)=\sigma(j)+2(c-1)\sigma(j/2)[2\mid j]$, which forms no $q$-product and
touches no pentagonal number, every division by $n$ being asserted exact ---
finds the two routes equal at all $2567$ coefficients, and compares the result
against the integers $W$ and $Q$ of \S\ref{sec:refute} \emph{as decimal
strings}, so a misprint in this paper would fail the program.

Two published sequences pin \eqref{eq:gf} against Remark~\ref{rem:defn}: at
$c=1$ the program reproduces OEIS A000041 and at $c=2$ OEIS A002513
\cite{OEIS}. The recorded run of \texttt{verify.py} reports
\begin{center}
\texttt{VERDICT: ALL 35 CHECKS PASS}
\end{center}
and the program exits nonzero if any check fails. The run reports more than this
note uses: residue-class sweeps, checks that congruences proved in
\cite{DasMaitySaikia} are not violated over the bounded ranges swept, and
residues for the other two displayed lines. None of that is claimed here, and
all of it is a statement about a bounded range. Its output also records what it
does not do: it proves none of the three congruences on any progression, it does
not compute the exact integers behind the other two lines, and it does not check
the source's Ramanujan--Kolberg parameter table against \cite{Radu} --- its final
section recomputes one set of residues from the formula as printed in the source,
and this note makes no claim about that table.

\section{Scope}\label{sec:scope}

Only \eqref{eq:line1}, as printed and quoted in \S\ref{sec:conj}, is refuted,
and only by exhibiting one $n$. Nothing here is a claim about $\vv_{31}$ at an
unlisted cell: the congruence does hold at $n=0,1,10$ of the window checked by
\texttt{verify.py}, with $\vv_{31}=2$ exactly at $n=0$ and $n=1$, and no
assertion is made about how often that happens.

Nothing new modulo the first power of $31$ is proved here: \eqref{eq:published}
is the $p=31$ case of Theorem 3 of \cite{Guadalupe} and is quoted, not proved.

No Ramanujan--Kolberg certificate is offered or completed, in either direction:
nothing checked here would have proved \eqref{eq:line1} had the cells examined
all vanished, and Theorem~\ref{thm:main} needs no certificate.

Finally, all quotations and parameter values are taken from the source of
\texttt{arXiv:\allowbreak 2503.19399v2}, submitted 28 April 2026, and the wording
quoted in \S\ref{sec:conj} is established here for that version only. The version
of record (\emph{Eur.\ J.\ Math.} \textbf{12} (2026), no.~3, Paper No.~34) was
revised on 26 May 2026 and accepted on 28 May 2026, that is, after that
\texttt{arXiv} version; its body is paywalled and was not read, so whether the
accepted version prints this conjecture in this wording is not established here.

\begin{thebibliography}{9}

\bibitem{ASS}
T.~Amdeberhan, J.~A.~Sellers, and A.~Singh,
\emph{Arithmetic properties for generalized cubic partitions and
overpartitions modulo a prime},
Aequationes Math. \textbf{99} (2025), 1197--1208.
\href{https://doi.org/10.1007/s00010-024-01116-7}{doi:10.1007/s00010-024-01116-7}.
The definition of $a_c(n)$ originates here, with the word \emph{even}.

\bibitem{DasMaitySaikia}
H.~Das, S.~Maity, and M.~P.~Saikia,
\emph{Arithmetic properties of generalized cubic and overcubic partitions},
\href{https://arxiv.org/abs/2503.19399}{arXiv:2503.19399v2} (28 April 2026);
published as Eur. J. Math. \textbf{12} (2026), no.~3, Paper No.~34, 22 pp.
The conjecture quoted here, and every quotation and parameter value, are from
the \texttt{arXiv} source of v2; the journal text was not consulted.

\bibitem{GhoshalJana}
S.~Ghoshal and A.~Jana,
\emph{Combinatorial proof of a result on generalized overcubic partitions and
related conjectures},
\href{https://arxiv.org/abs/2512.04775}{arXiv:2512.04775} (2025).
Quotes the third congruence of the same conjecture of \cite{DasMaitySaikia} as
``Conjecture 7.3''.

\bibitem{Guadalupe}
R.~Guadalupe,
\emph{A note on congruences for generalized cubic partitions modulo primes},
Integers \textbf{25} (2025), Article A20;
\href{https://math.colgate.edu/~integers/z20/z20.pdf}{available from the journal}.
Theorem 3 gives \eqref{eq:published} at $p=31$.

\bibitem{OEIS}
OEIS Foundation Inc.,
\emph{The On-Line Encyclopedia of Integer Sequences},
sequences \href{https://oeis.org/A000041}{A000041} ($a_1$, the partition
function) and \href{https://oeis.org/A002513}{A002513} ($a_2$, the cubic
partitions).

\bibitem{Radu}
S.~Radu,
\emph{An algorithmic approach to Ramanujan's congruences},
Ramanujan J. \textbf{20} (2009), 215--251.
\href{https://doi.org/10.1007/s11139-009-9174-0}{doi:10.1007/s11139-009-9174-0}.
Not obtained; cited only for the algorithm whose bound is reported in
\cite{DasMaitySaikia}.

\end{thebibliography}

\end{document}
